% modules/open-targets/A3-heavy-tailed-posteriors.tex
% -----------------------------------------------------------------------------
% Deep-dive companion to the A3 entry of the Part I agenda (modules/08-open-targets)
% and to Tier 4 (modules/06-tier4-heavy-tails): tail-index-dependent functional-
% inequality constants for heavy-tailed robust-Bayes posteriors -- weighted and
% weak Poincare inequalities, the exactly-solved Student-t/generalized-Cauchy model
% case, the sharp horseshoe theorem, and the hierarchical-posterior frontier.
% Compiles inside main.tex; standalone it shows cross-module \ref as "??".
% -----------------------------------------------------------------------------
\documentclass[../../main.tex]{subfiles}
\begin{document}
\section{A3 --- Heavy-tailed posteriors beyond classical LSI}
\label{sec:a3}

This is the deep-dive companion to the A3 paragraph of the agenda (Section~\ref{sec:agenda})
and the quantitative continuation of Tier~4 (Section~\ref{sec:tier4},
Family~\ref{fam:heavy-tail-classification}). The target is to replace the purely negative
Tier-4 verdict --- robust-Bayes priors (Student-$t$, Cauchy, horseshoe, scale mixtures of
Gaussians) satisfy no classical LSI and no $T_2$ --- by \emph{positive}, tail-index-dependent
constants for the inequalities that \emph{do} hold. The agenda slogan,
\begin{align*}
  &\text{Student-}t,\ \text{horseshoe},\ \text{scale mixtures} \\
  &\quad\Longrightarrow\quad
  \text{tail-index-dependent Poincar\'e constants via $1$D Hardy and tensorization,}
\end{align*}
is the right programme, but it must be sharpened on three points that this section makes
precise. First, the correct object is almost never the \emph{classical} Poincar\'e constant ---
for polynomial tails that fails too --- but a \emph{weighted} Poincar\'e constant with a
quadratically growing weight. Second, the Student-$t$/generalized-Cauchy model case is
\emph{not open}: the sharp optimal weighted spectral gaps are known for the whole parameter
range, and A3 should use them as the calibration theorem, not advertise them as a target.
Third, the horseshoe now has a sharp weighted Poincar\'e theorem, while the genuine frontier is the
general scale-mixture catalogue and, above all, propagation of these constants through a
likelihood and through global-local hierarchical dependence, where the posterior is no longer a
product. The honest framing is therefore: \emph{imported exactly solved model family $+$ sharp
horseshoe theorem $+$ a general scale-mixture / dependence programme.}

\noindent\textbf{Reviewed status.}
The sharp generalized-Cauchy result below is imported from the cited literature.  The horseshoe,
fixed-marginal obstruction, and block-Gibbs results have standalone solution dossiers and passed
independent agent review. The product and hierarchical-prior results below are likewise
independently agent-certified; the bounded-likelihood posterior transfer remains an analytic
draft.

\subsection{Why the right object is the weighted Poincar\'e constant}
\label{subsec:a3-weighted}

The first correction to the naive paragraph is structural: for polynomial tails the classical
Poincar\'e inequality is the wrong target, not merely the LSI. Theorem~\ref{thm:heavy-tail-no-lsi}
already rules out LSI and $T_2$ (an LSI forces sub-Gaussian concentration), but the classical
Poincar\'e inequality $\Var_\pi(f)\le\CP\int\norm{\nabla f}^2\dd\pi$ itself fails for the heaviest
tails. For the generalized Cauchy law
\begin{equation}\label{eq:a3-gencauchy}
  \dd\mu_\beta(x)\ \propto\ (1+\norm{x}^2)^{-\beta}\dd x,
  \qquad \beta>\tfrac d2,
\end{equation}
the $1$D Hardy criterion (Theorem~\ref{thm:hardy-1d}) diagnoses the failure cleanly: with
$p(x)\sim c\,\abs x^{-2\beta}$ in dimension one, the unweighted Hardy functional is
\[
  \mu_\beta([x,\infty))\int^x\frac{\dd t}{p(t)}
  \ \asymp\ x^{-(2\beta-1)}\,x^{2\beta+1}\ =\ x^{2}\ \to\ \infty,
\]
so $\CP(\mu_\beta)=\infty$. The repair is to weaken the Dirichlet form by a weight $a(x)$ that
grows in the tails. The \emph{weighted Poincar\'e inequality}
\begin{equation}\label{eq:a3-weighted-pi}
  \Var_\pi(f)\ \le\ C_w\int\inner{A(x)\nabla f(x)}{\nabla f(x)}\dd\pi(x),
  \qquad A(x)\ \asymp\ (1+\norm{x}^2)\,\Id,
\end{equation}
restores a finite constant, because the quadratic weight exactly cancels the polynomial growth in
the Hardy integral:
\[
  \mu_\beta([x,\infty))\int^x\frac{\dd t}{(1+t^2)p(t)}
  \ \asymp\ x^{-(2\beta-1)}\,x^{2\beta-1}\ =\ O(1).
\]
The correct replacement of the Tier-4 negative verdict is thus the hierarchy
\begin{equation}\label{eq:a3-hierarchy}
  \underbrace{\text{classical LSI}/T_2}_{\text{fail (Thm.~\ref{thm:heavy-tail-no-lsi})}}
  \ \leadsto\
  \underbrace{\text{weighted PI \eqref{eq:a3-weighted-pi}, weak PI, weighted LSI}}_{\text{hold, with tail-index-dependent constants}} .
\end{equation}
A weighted PI controls a \emph{preconditioned} diffusion whose carr\'e du champ is
$\Gamma(f)(x)=\inner{A(x)\nabla f}{\nabla f}$, not the vanilla overdamped Langevin diffusion;
the unweighted dynamics instead obeys a \emph{weak} Poincar\'e inequality
\begin{equation}\label{eq:a3-weak-pi}
  \Var_\pi(f)\ \le\ \beta_{\mathrm{WPI}}(s)\int\norm{\nabla f}^2\dd\pi\ +\ s\,\osc(f)^2,
  \qquad s>0,
\end{equation}
with a subgeometric rate; for a pure power tail of index $\alpha$ the truncation/Hardy heuristic
gives $\beta_{\mathrm{WPI}}(s)\asymp s^{-2/\alpha}$ \cite{CattiauxGozlanGuillinRoberto2010}.
For an iid $d$-fold product the natural tensorization gives
$\beta_d(s)\le\beta_1(s/d)$, hence a $d^{2/\alpha}$ loss; only the weighted metric below is
dimension-free.
Which object A3 wants depends on the sampler: \eqref{eq:a3-weighted-pi} for a Riemannian /
preconditioned scheme, \eqref{eq:a3-weak-pi} for plain Langevin.

\begin{warning}[A light-tailed likelihood can restore the classical inequalities]
\label{warn:a3-likelihood}
The whole tier is about posteriors that \emph{inherit} polynomial tails. A Student-$t$ prior
combined with a full-rank Gaussian likelihood yields a posterior with Gaussian tails, where
classical LSI and Poincar\'e return and the prior tail index is irrelevant. A3 is therefore
genuinely about the regimes where the tail survives into the posterior: weak identification,
bounded or saturating likelihoods, underdetermined directions, and the global-/local-scale
\emph{hyperparameter} marginals of hierarchical priors, which are heavy-tailed by construction.
This is the analogue of the flat-direction obstruction of A1
(Proposition~\ref{prop:flat-prior-nonintegrable}), read in the tail rather than the bulk.
\end{warning}

\subsection{The \texorpdfstring{$1$D}{1D} Hardy certificate is exactly the right tool}
\label{subsec:a3-hardy}

The engine for every one-dimensional constant in this tier is the weighted form of the Hardy /
Muckenhoupt criterion (Theorem~\ref{thm:hardy-1d}). For $\mu(\dd x)=p(x)\dd x$ on $\R$ with median
$m$ and positive weight $a(x)$, the weighted Poincar\'e inequality
$\Var_\mu(f)\le C\int a\,\abs{f'}^2\dd\mu$ holds iff the two Hardy quantities
\begin{equation}\label{eq:a3-hardy-B}
  B_+\ =\ \sup_{x>m}\ \mu([x,\infty))\int_m^x\frac{\dd t}{a(t)\,p(t)},
  \qquad
  B_-\ =\ \sup_{x<m}\ \mu((-\infty,x])\int_x^m\frac{\dd t}{a(t)\,p(t)},
\end{equation}
are finite, and then the optimal constant is pinned up to a universal factor $4$,
\begin{equation}\label{eq:a3-hardy-factor4}
  \max(B_+,B_-)\ \le\ C_{\mathrm{opt}}\ \le\ 4\max(B_+,B_-).
\end{equation}
This is the weighted reading of Theorem~\ref{thm:hardy-1d} (the stated criterion is the case
$a\equiv1$); Miclo's analysis \cite{Miclo2008Hardy} characterizes when such Hardy bounds even
pin the \emph{exact} one-dimensional constant, not just up to factor $4$. The mechanism of
\eqref{eq:a3-hardy-B}--\eqref{eq:a3-hardy-factor4} is the entire reason A3 is tractable in one
dimension: it turns ``does a tail-index-dependent constant exist, and what is it'' into a single
explicit supremum, evaluable in closed form for the model families below and numerically for the
rest. Combined with tensorization (Subsection~\ref{subsec:a3-tensor}) it is the whole toolkit for
product priors.

\subsection{The \texorpdfstring{Student-$t$}{Student-t}/generalized-Cauchy model case is solved --- use it to calibrate}
\label{subsec:a3-student}

Here the naive paragraph most overstates the openness. For the generalized Cauchy family
\eqref{eq:a3-gencauchy} the sharp optimal weighted spectral gap is known for the entire range
$\beta>d/2$ \cite{Huguet2024Cauchy}, building on the one-dimensional diffusion viewpoint of
\cite{BonnefontJoulinMa2016}, who already noted that the classical operator has no spectral gap
while the weight $1+x^2$ repairs it.

\begin{theorem}[Sharp weighted Poincar\'e constant for generalized Cauchy / Student-$t$; {\rm\cite{Huguet2024Cauchy,BonnefontJoulinMa2016}}]
\label{thm:a3-student}
For $\mu_\beta$ as in \eqref{eq:a3-gencauchy} with $\beta>d/2$, the sharp inequality
\[
  \lambda_{\beta,d}\,\Var_{\mu_\beta}(f)\ \le\ \int_{\R^d}(1+\norm{x}^2)\,\norm{\nabla f}^2\dd\mu_\beta
\]
holds with the explicit optimal spectral gap
\[
  \lambda_{\beta,1}=
  \begin{cases}
    (\beta-\tfrac12)^2, & \tfrac12<\beta\le\tfrac32,\\[1mm]
    2(\beta-1), & \beta\ge\tfrac32,
  \end{cases}
  \qquad
  \lambda_{\beta,d}=
  \begin{cases}
    (\beta-\tfrac d2)^2, & \tfrac d2<\beta\le\tfrac d2+2,\\[1mm]
    4(\beta-\tfrac d2-1), & \tfrac d2+2\le\beta\le d+1,\\[1mm]
    2(\beta-1), & \beta\ge d+1.
  \end{cases}
\]
Equivalently $\CP^{w}(\mu_\beta)=\lambda_{\beta,d}^{-1}$ for the weight $A=(1+\norm{x}^2)\Id$.
\end{theorem}

Translating to the standard $d$-dimensional Student-$t_\nu$ law
$\pi_{\nu,d}\propto(1+\norm{x}^2/\nu)^{-(\nu+d)/2}$ via $\beta=(\nu+d)/2$ and the rescaling
$x=\sqrt\nu\,y$ gives the sharp weighted inequality with tail weight $\nu+\norm{x}^2$; in one
dimension, with the normalized weight $1+x^2/\nu$,
\begin{equation}\label{eq:a3-student-1d}
  \Var_{t_\nu}(f)\ \le\ C_{\nu}^{\mathrm{norm}}\int\Big(1+\tfrac{x^2}{\nu}\Big)\abs{f'}^2\dd t_\nu,
  \qquad
  C_\nu^{\mathrm{norm}}=
  \begin{cases}
    4/\nu, & 0<\nu\le2,\\[1mm]
    \nu/(\nu-1), & \nu\ge2.
  \end{cases}
\end{equation}
The constant degrades exactly as the tail heavies ($C_\nu^{\mathrm{norm}}\to\infty$ as $\nu\to0$,
i.e.\ tail index $\alpha=\nu\to0$). The weight normalization must be tracked: for the unscaled
generalized-Cauchy density and weight $1+x^2$, the exact one-dimensional branches are
$C_w=4/\alpha^2$ for $0<\alpha\le2$ and $C_w=1/(\alpha-1)$ for $\alpha\ge2$. Thus
$\alpha^{-2}$ is the near-integrability branch, not a uniform law over all tail indices.
Huguet's constant is the \emph{radial} one for the elliptically contoured multivariate law; for an independent product of
Student-$t$ coordinates one instead tensorizes the $1$D constant \eqref{eq:a3-student-1d}
(Subsection~\ref{subsec:a3-tensor}), which is simpler and dimension-free in the weighted metric.

\begin{remark}[Bobkov--Ledoux: the convex backbone]
\label{rem:a3-bobkovledoux}
The Student-$t$ story is not isolated. For densities $p=V^{-\beta}$ with $V$ positive and convex
and $\beta>n$, Bobkov--Ledoux \cite{BobkovLedoux2009} prove a weighted Poincar\'e inequality with
a universal quadratic weight of the form $r^2+\norm{x}^2$ ($r$ a fixed quantile radius), and a
weighted log-Sobolev inequality with the stronger weight $(1+\norm{x}^2)^2$,
$\Ent_{\nu_\beta}(g^2)\le\frac1{\beta-1}\int\norm{\nabla g}^2(1+\norm{x}^2)^2\dd\nu_\beta$. The
structural principle they encode --- \emph{polynomial tail $\Rightarrow$ quadratic tail weight in
the Dirichlet form} --- is the backbone A3 should cite for any convex / $\kappa$-concave target;
the weighted PI is the sharper and more tail-sensitive object, the weighted LSI the fallback when
entropy control is genuinely needed.
\end{remark}

\subsection{Horseshoe: sharp weighted constant}
\label{subsec:a3-horseshoe}

The horseshoe prior \cite{CarvalhoPolsonScott2010,CarvalhoPolsonScott2009} is the canonical
normal scale mixture, $\theta_i\mid\lambda_i,\tau\sim N(0,\lambda_i^2\tau^2)$ and
$\lambda_i\sim C^+(0,1)$. Its one-dimensional marginal has the exact special-function form
\begin{equation}\label{eq:a3-horseshoe-density}
  h_\tau(\theta)
  =\frac{e^aE_1(a)}{\pi\sqrt{2\pi}\,\tau},
  \qquad a=\frac{\theta^2}{2\tau^2},
  \qquad E_1(a)=\int_a^\infty\frac{e^{-s}}s\,\dd s .
\end{equation}
It has a logarithmic pole at the origin and $h_\tau(\theta)\asymp\tau/\theta^2$ in the tails
\cite{CarvalhoPolsonScott2010,Castillo2017HorseshoeDiscussion}. The Cauchy-type tail makes the
classical Poincar\'e inequality fail by the calculation of Subsection~\ref{subsec:a3-weighted}, but
the Cauchy weight $a_\tau(\theta)=\tau^2+\theta^2$ is essentially forced and the resulting weighted
inequality is bounded through the Hardy criterion.  The Hardy bracket gives a computable first
certificate; the intrinsic-coordinate argument below identifies the exact constant.

\begin{proposition}[Sharp weighted Poincar\'e inequality for the horseshoe marginal]
\label{prop:a3-horseshoe}
Let $\mathrm{HS}_\tau$ have density \eqref{eq:a3-horseshoe-density} and define, by symmetry,
\[
  B_{\mathrm{HS}}(\tau)\ =\ \sup_{r>0}\ \Prob_{\mathrm{HS},\tau}(\theta\ge r)
  \int_0^r\frac{\dd t}{(\tau^2+t^2)\,h_\tau(t)} .
\]
Then $B_{\mathrm{HS}}(\tau)<\infty$, and the weighted Poincar\'e inequality
\[
  \Var_{\mathrm{HS},\tau}(f)\ \le\ C_{\mathrm{HS}}(\tau)\int(\tau^2+\theta^2)\,\abs{f'(\theta)}^2\dd\mathrm{HS}_\tau(\theta)
\]
holds with the sharp constant
\[
  B_{\mathrm{HS}}(\tau)=B_{\mathrm{HS}}(1),\qquad
  C_{\mathrm{HS}}(\tau)=C_{\mathrm{HS}}(1)=4.
\]
\end{proposition}

\begin{proof}
The Hardy finiteness follows directly from \eqref{eq:a3-horseshoe-density}: the logarithmic pole
makes $1/h_\tau$ locally integrable at the origin, while in the tail
$\Prob(\theta\ge r)\asymp\tau/r$ and
$\int^r\dd t/((\tau^2+t^2)h_\tau(t))\asymp r/\tau$.  The substitution
$\theta=\tau x$ gives both scale identities, so it remains to prove the sharp value at $\tau=1$.

Put
\[
  H(a)=e^aE_1(a)=\int_0^\infty\frac{e^{-t}}{a+t}\,\dd t,
  \qquad a>0.
\]
The fourth Laplace continued-fraction convergent has a global upper remainder of the needed sign:
\begin{equation}\label{eq:a3-e1-upper}
 H(a)<\frac{a^2+5a+2}{a(a^2+6a+6)}.
\end{equation}
Indeed, with $N=a^2+5a+2$ and $Q=a(a^2+6a+6)$, polynomial division followed by two integrations
by parts gives the exact remainder identity
\[
 N-\frac{Q}{a+t}
 =a(t-1)-t^2+6t-4+\frac{t(t^2-6t+6)}{a+t}.
\]
The first two terms integrate to zero against $e^{-t}\dd t$, since
$\int e^{-t}t\dd t=1$ and $\int e^{-t}t^2\dd t=2$.  Therefore
\[
 N-QH(a)
 =\int_0^\infty\frac{\{e^{-t}t^3\}''}{a+t}\,\dd t
 =2\int_0^\infty\frac{e^{-t}t^3}{(a+t)^3}\,\dd t>0.
\]
For the two integrations by parts, $e^{-t}t^3=(e^{-t}t^3)'=0$ at $t=0$, and both quantities
converge to zero at infinity; hence every boundary term vanishes.

Now use the intrinsic coordinate $y=\operatorname{arsinh}x$.  The weighted Dirichlet form becomes
the ordinary form $\int\abs{f'(y)}^2q(y)\dd y$, where, up to normalization,
\[
 q(y)=\cosh y\,H\!\left(\frac{\sinh^2y}{2}\right),
 \qquad W(y)=-\log q(y).
\]
For $y>0$, set $x=\sinh y$, $c=\cosh y$, and $a=x^2/2$.  Since
$H'(a)=H(a)-1/a$,
\[
 W'(y)=xc\left(\frac1{aH(a)}-1\right)-\frac{x}{c}.
\]
Inequality \eqref{eq:a3-e1-upper} implies
\[
 \frac1{aH(a)}\ge\frac{a^2+6a+6}{a^2+5a+2}
 \ge 1+\frac1{c^2}+\frac1{c(c+1)}.
\]
The last inequality is purely algebraic; more precisely,
\[
 \frac{a^2+6a+6}{a^2+5a+2}
 -\left(1+\frac1{c^2}+\frac1{c(c+1)}\right)
 =\frac{c^2(c-1)^2+5c^2+2c+1}
 {c^2(c+1)(c^4+8c^2-1)}>0.
\]
Consequently
\begin{equation}\label{eq:a3-horseshoe-drift}
 W'(y)-\tanh(y/2)
 \ge\frac{x\{c^2(c-1)^2+5c^2+2c+1\}}
 {c(c+1)(c^4+8c^2-1)}>0,
 \qquad y>0.
\end{equation}

Let $\mathcal Lf=f''-W'f'$ be the intrinsic diffusion generator on $(0,\infty)$ and take
$g(y)=\sinh(y/2)$.  Then
\[
 -\frac{\mathcal Lg}{g}
 =-\frac14+\frac12W'(y)\coth(y/2)\ge\frac14.
\]
For a compactly supported Dirichlet test $\varphi=gu$, the ground-state identity therefore gives
\[
 \int_0^\infty\!\left(\abs{\varphi'}^2-\frac14\varphi^2\right)q\,\dd y
 =\int_0^\infty g^2\abs{u'}^2q\,\dd y
  +\int_0^\infty\left(-\frac{\mathcal Lg}{g}-\frac14\right)\varphi^2q\,\dd y
 \ge0.
\]
Let $\mathsf D(\mathcal E)$ be the closure of $C_c^\infty(\R)$ for
$\norm{f}_{L^2(q)}^2+\int\abs{f'}^2q$.  Since $q(y)\asymp\log(1/y)$ at $0$, both
$q$ and $1/q$ are locally integrable.  Elements of the form domain therefore have locally
absolutely continuous representatives and a common two-sided trace at the origin, with
\[
 \abs{f(y)-f(0)}^2
 \le\left(\int_0^y\abs{f'}^2q\right)
     \left(\int_0^y\frac{\dd s}{q(s)}\right).
\]
Odd functions have trace zero.  Symmetrizing the defining $C_c^\infty(\R)$ core shows that smooth
compactly supported odd functions are a core for the odd domain.  For such a test
$u=\varphi/g=O(1)$ at zero, and the boundary term in the ground-state calculation is
$qgg'u^2=O(y\log(1/y))\to0$.  Identity (and hence the inequality) above holds first on this odd
core; form closure extends the inequality, without asserting convergence of the two separate
nonnegative terms on the right, to the entire odd domain.
By odd reflection, every trace-zero half-line form function therefore has Rayleigh quotient at
least $1/4$.

This Dirichlet estimate controls the even sector as well, without a Sturm interlacing argument.
Indeed, let $f$ be an even-sector half-line form function with
$\int_0^\infty f q\,\dd y=0$ and set $h=f-f(0)$.  Then $h$ has trace zero,
$\mathcal E(h)=\mathcal E(f)$, and
\[
 \int_0^\infty h^2q\,\dd y
 =\int_0^\infty f^2q\,\dd y
   +f(0)^2\int_0^\infty q\,\dd y
 \ge\int_0^\infty f^2q\,\dd y.
\]
The Dirichlet estimate applied to $h$ therefore bounds every mean-zero even function by the same
$1/4$ gap.  Parity decomposition is orthogonal in both $L^2(q)$ and the Dirichlet form, so the
full gap is at least $1/4$.

For the reverse inequality, the tail expansion
$H(a)=a^{-1}-a^{-2}+O(a^{-3})$ yields
$q(y)=K e^{-y}(1+O(e^{-2y}))$.  Take increasingly wide smooth cutoffs of $e^{y/2}$ supported in
$[R,R+L]$.  Direct substitution in the Rayleigh quotient gives
$1/4+O(L^{-2})+O(e^{-2R})$; subtracting the $q$-mean changes no energy and a negligible fraction
of the squared norm.  Sending first $R$ and then $L$ to infinity shows that the full gap is at
most $1/4$.  Therefore it is exactly $1/4$, i.e. $C_{\mathrm{HS}}=4$.
\end{proof}

\subsection{General scale-mixture priors: a Hardy-computable certificate, not exact constants}
\label{subsec:a3-scalemix}

Most robust-Bayes priors are normal scale mixtures $\theta\mid s\sim N(0,s^2)$, $s\sim\rho$ ---
horseshoe, Student-$t$, Bayesian lasso/Laplace, normal-gamma, horseshoe$+$, and the global-local
shrinkage family, which is regularly varying \cite{BhadraDattaPolsonWillard2016}. The natural
classification is by marginal tail index: if $p(\theta)\sim c\,\abs\theta^{-1-\alpha}L(\abs\theta)$
with $L$ slowly varying and $\alpha>0$, then classical PI fails, LSI and $T_2$ fail more strongly.
For $a(\theta)\sim\kappa\theta^2$, Karamata's theorem gives the exact \emph{tail limit}
\begin{equation}\label{eq:a3-alpha-scaling}
  \mu([x,\infty))\int^x\frac{\dd t}{a(t)p(t)}\longrightarrow
  \frac{1}{\kappa\alpha^2}.
\end{equation}
This forces the Hardy tail floor but does not upper-bound its full supremum: a finite-location
bottleneck can dominate. Consequently $C_w\asymp\alpha^{-2}$ as $\alpha\downarrow0$ requires
uniform $O(\alpha^{-2})$ control of the bulk Hardy contribution. This matches the heaviest-tail
branch of Theorem~\ref{thm:a3-student}, but not its lighter-tail branches. For the broad
scale-mixture class A3 should therefore promise the
Hardy-computable certificate
\begin{equation}\label{eq:a3-certificate}
  C_w\ \approx\ \sup_x\ \mu([x,\infty))\int^x\frac{\dd t}{(1+t^2)\,p(t)},
\end{equation}
with the tail floor \eqref{eq:a3-alpha-scaling}, an explicit bulk check, and numerical evaluation
when a closed form is unavailable.

\begin{remark}[Stein-kernel route: useful, but not at the Cauchy boundary]
\label{rem:a3-stein}
There is an elegant alternative in one dimension: every measure with finite first moment and
connected support satisfies a weighted Poincar\'e inequality with its \emph{Stein kernel} as
weight \cite{Saumard2019SteinKernel}. This is attractive for scale mixtures with tail index
$\alpha>1$, where a first moment exists. It is \emph{not} the right primary tool at the standard
horseshoe / Cauchy boundary $\alpha=1$, where the absolute first moment diverges; there the Hardy
route of Subsection~\ref{subsec:a3-hardy} is the robust one.
\end{remark}

\subsection{Tensorization: product priors are dimension-free, hierarchical posteriors are not}
\label{subsec:a3-tensor}

\begin{theorem}[Product Hardy certificate]
\label{thm:a3-product}
Let $d\ge1$ be finite and $\mu=\bigotimes_{i=1}^d\mu_i$. If each coordinate satisfies
$\Var_{\mu_i}(g)\le C_i\int a_i\abs{g'}^2\dd\mu_i$, then
\begin{equation}\label{eq:a3-tensor}
  \Var_\mu(f)\ \le\ \Big(\max_i C_i\Big)\int\sum_i a_i(x_i)\,\abs{\partial_i f}^2\dd\mu.
\end{equation}
In particular, if $B_i=\max(B_{i,+},B_{i,-})$ denotes the two-sided Hardy quantity of
$(\mu_i,a_i)$, the constant on the right is at most $4\max_iB_i$.
\end{theorem}

\begin{proof}
Conditional-variance tensorization gives \eqref{eq:a3-tensor}; the weighted Hardy criterion gives
$C_i\le4B_i$ in each coordinate.
\end{proof}

This theorem gives a \emph{dimension-free} weighted constant for independent Student-$t$,
independent horseshoe marginals (weight $\tau^2+\theta_j^2$), and independent regularly varying
scale mixtures. It is the independently reviewed easy half of the tensorization clause.

The hard half is the genuine \emph{hierarchical posterior}. With a global scale $\tau$, local
scales $\lambda_j$, and a likelihood $L(y\mid\theta)$,
\[
  \pi(\dd\theta,\dd\lambda,\dd\tau\mid y)\ \propto\ L(y\mid\theta)
  \prod_j N(\dd\theta_j;0,\tau^2\lambda_j^2)\,p(\dd\lambda_j)\,p(\dd\tau),
\]
which is \emph{not} a product. The conditional-variance decomposition
\[
  \Var_\pi(f)\ =\ \E\big[\Var(f\mid\lambda,\tau)\big]\ +\ \Var\big(\E[f\mid\lambda,\tau]\big)
\]
has conditional Gaussian Poincar\'e constants scaling like $\tau^2\lambda_j^2$, which are
\emph{unbounded} as the scales wander; one should therefore expect a weighted \emph{joint}
inequality, not a Euclidean one. Combining conditional Gaussian Poincar\'e with Hardy inequalities
for $(\lambda,\tau)$ and a two-scale / decomposition argument to obtain posterior-specific weighted
constants is the genuinely open Bayesian direction, and is especially delicate when the likelihood
couples coordinates or changes the tail behaviour (Warning~\ref{warn:a3-likelihood}).

There is also a concrete certified prior-side result in the natural non-centred metric.

\begin{proposition}[Hierarchical horseshoe prior in the pullback metric]
\label{prop:a3-hierarchical-prior}
Fix a finite $d\ge1$. Let $z_1,\ldots,z_d\sim N(0,1)$ and let
$u,v_1,\ldots,v_d$ be independent log-half-Cauchy variables, all mutually independent, and set
$\theta_j=e^{u+v_j}z_j$ for $1\le j\le d$. The joint prior satisfies a weighted Poincar\'e inequality with
optimal constant $4$ and Dirichlet energy
\[
 \int\!\left\{
 \sum_j e^{2u+2v_j}\abs{\partial_{\theta_j}f}^2
 +\sum_j\abs{\partial_{v_j}f+\theta_j\partial_{\theta_j}f}^2
 +\abs{\partial_uf+\sum_j\theta_j\partial_{\theta_j}f}^2
 \right\}\dd\pi .
\]
\end{proposition}

\begin{proof}
A log-half-Cauchy variable has density $(\pi\cosh u)^{-1}$ and classical Poincar\'e constant
$4$, by transforming the sharp weighted Cauchy inequality of
Theorem~\ref{thm:a3-student} through $u=\operatorname{arsinh}x$. The base product law of
$(z,u,v)$ therefore has constant $4$. Tensorization and the chain rule under
$\theta_j=e^{u+v_j}z_j$ give the displayed pullback energy. Testing a function of $u$ alone shows
that the constant is optimal.
\end{proof}

\noindent Both results are independently agent-certified in the shared standalone dossier
\begin{center}\small
\path{solutions/a3-product-hierarchical-prior.tex}.
\end{center}
The review checks the sharp tensorization constant, log-half-Cauchy normalization, complete
pullback cross terms, and optimality.

\subsection{Dependence obstruction and replacement target}
\label{subsec:a3-theorem}

\begin{warning}[Retired marginal-only conjecture]
\label{warn:a3-marginal-only}
The former A3 conjecture asserted that the joint weighted Poincar\'e constant of a heavy-tailed
posterior could be controlled only by one-dimensional Hardy quantities of its coordinate
marginals. This is false: fixed marginals do not control a dependence bottleneck. The certified
analytic counterexample below records the obstruction; the viable replacement must include a
quantitative dependence parameter, as in Theorem~\ref{thm:a3-block-gibbs} and
Conjecture~\ref{conj:a3-dependent}.
\end{warning}

\begin{proposition}[Fixed Cauchy marginals do not control the joint gap]
\label{prop:a3-marginals-not-joint}
For every $0<\eps<1$ and in every dimension at least three, there are positive, piecewise-smooth
joint densities whose coordinate marginals are all standard Cauchy (and hence have the same fixed
finite Hardy quantities for the weight $1+x^2$), but whose joint weighted Poincar\'e constants in
the diagonal metric $\diag(1+x_1^2,\ldots,1+x_d^2)$ are of order $\eps^{-1}$.
\end{proposition}

\begin{proof}
It suffices to work in $d=3$; higher dimensions follow by adjoining independent Cauchy
coordinates.  Let $p(x)=\{\pi(1+x^2)\}^{-1}$ have CDF $F$, put
$A=(0,1/2)$ and $D=(1/2,1)$, and on $(0,1)^3$ use the positive copula density
\[
 c_\eps=(1-\eps)4(\mathbf1_{A^3}+\mathbf1_{D^3})+\eps.
\]
Every one-dimensional marginal of $c_\eps$ is uniform, so
$q_\eps(x)=\prod_jp(x_j)c_\eps(F(x_1),F(x_2),F(x_3))$ has marginal $p$ in every coordinate.

To construct the bottleneck test, translate the common corner
$o=(1/2,1/2,1/2)$ to the origin.  Choose a smooth odd function $g$ on $S^2$ that equals $1$ on
the all-negative spherical octant, equals $-1$ on the all-positive octant, and changes only
through the six mixed-sign octants.  Then
$\varphi(u)=g((u-o)/\abs{u-o})$ (with an arbitrary value at $o$) belongs to
$H^1((0,1)^3)$: near $o$, $\abs{\nabla\varphi}=O(\abs{u-o}^{-1})$, whose square is integrable in
dimension three.  Moreover $\varphi=1$ on $A^3$, $\varphi=-1$ on $D^3$, and its gradient is
supported in the mixed octants, where $c_\eps=\eps$.  Symmetry gives mean zero, while the two
constant boxes have total mass $1-3\eps/4$, so
$\Var_{c_\eps}(\varphi)\ge1-3\eps/4$.

Finally set $f(x)=\varphi(F(x_1),F(x_2),F(x_3))$.  Under $u_j=F(x_j)$,
\[
 \int\sum_{j=1}^3(1+x_j^2)\abs{\partial_jf}^2q_\eps(x)\,\dd x
 =\int c_\eps(u)\sum_{j=1}^3
   (1+F^{-1}(u_j)^2)p(F^{-1}(u_j))^2\abs{\partial_j\varphi}^2\,\dd u
 \le \frac{\eps}{\pi^2}\int\norm{\nabla\varphi}^2\dd u.
\]
Thus the optimal joint constant is at least $c/\eps$. Conversely,
$\eps\le c_\eps\le4-3\eps$, so density comparison with the product Cauchy law gives an upper
bound $4(4-3\eps)/\eps$. More explicitly, smooth $H^1$ approximants
to $\varphi$ converge both in variance (because $c_\eps$ is bounded) and in the displayed energy
(because its coefficient is bounded by $4/\pi^2$).  For each fixed $\eps$, choosing the
approximation error $o(\eps)$ gives the same Rayleigh lower bound with smooth test functions.
No smoothing of the copula is required: $q_\eps$ is positive almost everywhere, and the
proposition does not assert smoothness.
\end{proof}

\noindent This obstruction is independently agent-certified in
\path{solutions/obs-marginals-not-joint.tex}; the audit scope is recorded in
\path{research/reviews/2026-08-22-a3-marginals-independent-audit.md}.

The dependence loss can be isolated in a second, purely probabilistic spectral gap.

\begin{theorem}[Block-Gibbs tensorization]
\label{thm:a3-block-gibbs}
Let $X=(X_1,\ldots,X_m)\sim\pi$ and define
\begin{equation}\label{eq:a3-block-gap}
 \gamma_{\rm blk}(\pi)
 :=\inf_{\Var_\pi(f)>0}
 \frac{\sum_{i=1}^m\E_\pi[\Var(f\mid X_{-i})]}{\Var_\pi(f)}.
\end{equation}
Suppose that, for almost every $x_{-i}$, the conditional law of $X_i$ satisfies
\[
 \Var_{\pi_i(\cdot\mid x_{-i})}(g)
 \le C_i\int\inner{A_i(x)\nabla_i g}{\nabla_i g}
          \dd\pi_i(\cdot\mid x_{-i}),
\]
with deterministic $C_i<\infty$.  If $\gamma_{\rm blk}(\pi)>0$, then
\begin{equation}\label{eq:a3-block-tensorization}
 \Var_\pi(f)
 \le \frac{\max_iC_i}{\gamma_{\rm blk}(\pi)}
 \int\sum_i\inner{A_i(x)\nabla_i f}{\nabla_i f}\dd\pi.
\end{equation}
For a product law $\gamma_{\rm blk}=1$.  For exactly two blocks $X,Y$,
\begin{equation}\label{eq:a3-hgr-gap}
 \gamma_{\rm blk}=1-\rho_{\max}(X,Y),
\end{equation}
where $\rho_{\max}$ is the Hirschfeld--Gebelein--R\'enyi maximal correlation; the equality is
understood in the extended sense, so a nonconstant common factor gives both
$\rho_{\max}=1$ and $\gamma_{\rm blk}=0$.
\end{theorem}

\begin{proof}
The definition of $\gamma_{\rm blk}$ followed by the conditional inequalities gives
\eqref{eq:a3-block-tensorization}.  For a product, the Efron--Stein/tensorization inequality gives
$\gamma_{\rm blk}\ge1$, while a function of one nondegenerate block gives the reverse inequality.

For two blocks, work in the centred Hilbert space $L^2_0(\pi)$ and let $P_X,P_Y$ be conditional
expectation onto the $X$- and $Y$-measurable subspaces.  Then
\[
 \E\Var(f\mid X)+\E\Var(f\mid Y)
 =\inner{f}{(2I-P_X-P_Y)f}.
\]
The norm of $P_XP_Y$ from the centred $Y$-subspace to the centred $X$-subspace is precisely
$\rho_{\max}$.  The two-projections identity
$\norm{P_X+P_Y}=1+\rho_{\max}$ on $L^2_0(\pi)$ therefore yields
\[
 \inf\operatorname{spec}(2I-P_X-P_Y)=1-\rho_{\max},
\]
which is exactly \eqref{eq:a3-hgr-gap}.  This also covers non-attainment by taking infima and
operator norms.
\end{proof}

An immediate posterior rung is available in the pullback metric of
Proposition~\ref{prop:a3-hierarchical-prior}.  If $\pi_0$ is that prior and
$\dd\pi_L=Z^{-1}L\,\dd\pi_0$ with $0<m\le L\le M<\infty$, density comparison gives
\begin{equation}\label{eq:a3-bounded-likelihood}
 C_w(\pi_L)\le4\frac{M}{m}
\end{equation}
for the same pullback Dirichlet form.  Indeed
$\Var_{\pi_L}(f)\le(M/Z)\Var_{\pi_0}(f)$ and the prior energy is at most
$(Z/m)$ times the posterior energy.  This bounded-oscillation corollary is deliberately only the
first rung; most Gaussian and logistic likelihoods have $\inf L=0$.

The next concrete model is the factorized horseshoe normal-means posterior
\[
 y_j\mid\theta_j\sim N(\theta_j,\sigma_j^2),\qquad
 \theta_j=e^{u+v_j}z_j.
\]
After integrating out $z_j$, its scale likelihood is, up to a constant,
\[
 m_j(w)=(\sigma_j^2+e^{2w})^{-1/2}
 \exp\!\left\{-\frac{y_j^2}{2(\sigma_j^2+e^{2w})}\right\},
 \qquad w=u+v_j,
\]
and the $(u,v)$ posterior is proportional to
$\operatorname{sech}(u)\prod_j\operatorname{sech}(v_j)m_j(u+v_j)$.  Conditional on $u$, the local blocks factorize.  A
model-specific proof of Conjecture~\ref{conj:a3-dependent} should therefore establish uniform
one-dimensional Hardy bounds for these conditional scale laws and a positive block gap
\eqref{eq:a3-block-gap} (or an explicit two-scale score-covariance bound).  Any loss in dimension,
data size, or standardized observations must remain visible through that coupling factor.

\begin{conjecture}[Dependence-aware hierarchical posterior inequality]
\label{conj:a3-dependent}
For a specified global--local hierarchy, conditional weighted Poincar\'e inequalities for the
coefficient and scale blocks, together with a quantitative inter-block coupling bound, imply a
joint weighted inequality in the non-centred pullback metric. The bound must depend explicitly on
the coupling parameter and must include derivatives in the scale variables and the structural
cross terms of Proposition~\ref{prop:a3-hierarchical-prior}; marginal tail indices alone are not
sufficient.
\end{conjecture}

\subsection{Sub-targets and what remains open}
\label{subsec:a3-open}

The imported generalized-Cauchy calibration and the independently reviewed sharp horseshoe and
block-Gibbs statements are already incorporated above. Together with the explicitly unreviewed
product and bounded-likelihood drafts, they leave the following research questions.

\begin{question}[Likelihood-tilted horseshoe constants]
\label{q:a3-horseshoe}
Proposition~\ref{prop:a3-horseshoe} settles the unperturbed horseshoe marginal.  For a
likelihood-tilted marginal $\dd\pi_L\propto L(\theta)h_\tau(\theta)\dd\theta$, determine
conditions weaker than bounded density ratio that preserve a finite weighted gap in the metric
$\tau^2+\theta^2$, quantify its departure from the prior value $4$, and identify when a
light-tailed likelihood instead changes the appropriate tail metric.  The first target is the
one-dimensional Gaussian normal-means likelihood, followed by logistic-type tilts.
\end{question}

\begin{question}[The scale-mixture catalogue]
\label{q:a3-catalogue}
Turn \eqref{eq:a3-certificate} into a catalogue: for each robust-Bayes prior (Laplace, normal-gamma,
horseshoe$+$, general regularly varying global-local mixtures \cite{BhadraDattaPolsonWillard2016})
record the tail index, the weighted Hardy constant (closed form or numeric), and the Stein-kernel
constant where $\alpha>1$ \cite{Saumard2019SteinKernel}. The deliverable is the Bayesian-prior
constant table that the general functional-inequality literature does not provide.
\end{question}

\begin{question}[Weak Poincar\'e rates for unweighted dynamics]
\label{q:a3-weak}
For samplers that see the plain Euclidean Dirichlet form, establish the weak Poincar\'e inequality
\eqref{eq:a3-weak-pi} with the tail-index rate $\beta_{\mathrm{WPI}}(s)\asymp s^{-2/\alpha}$ for the
posterior coordinate marginals, via the heavy-tailed functional-inequality toolkit
\cite{CattiauxGozlanGuillinRoberto2010,CattiauxGuillin2010}, and convert it into a subgeometric
convergence rate for Langevin on the raw parameterization. Record explicitly the
$d^{2/\alpha}$ product degradation rather than importing the dimension-free weighted conclusion.
\end{question}

\begin{question}[Hierarchical posteriors and likelihood perturbation]
\label{q:a3-hierarchical}
The genuinely open Bayesian part is Conjecture~\ref{conj:a3-dependent}: propagate conditional
weighted constants through a likelihood and hierarchical dependence.  First implement the
explicit block contract \eqref{eq:a3-block-tensorization} for the factorized horseshoe normal-means
model above; if the block gap degenerates, use a quantified two-scale score bound instead. The
metric must include scale derivatives and the cross terms exposed by
Proposition~\ref{prop:a3-hierarchical-prior}; Proposition~\ref{prop:a3-marginals-not-joint}
rules out any marginal-only theorem. Keep all dimension/data losses visible, then combine Hardy
control of the scale blocks with the perturbed-measure machinery of
\cite{CattiauxGuillin2022Perturbed} for less factorized likelihoods.
\end{question}

\begin{remark}[Relation to the rest of the agenda]
\label{rem:a3-relations}
A3 is the heavy-tailed face of the agenda. It shares the flat-direction / tail obstruction of the
finite-sample flagship A1 (Section~\ref{sec:a1}, Proposition~\ref{prop:flat-prior-nonintegrable}),
but in the tail rather than the bulk, and where A1's likelihood curvature \emph{repairs} the bulk,
here a light-tailed likelihood is what can repair the tail (Warning~\ref{warn:a3-likelihood}).
Asymptotically, a polynomial-tailed finite-$n$ posterior need not invalidate a classical
total-variation Bernstein--von Mises theorem for its local Gaussian core.  Total variation alone,
however, does not transfer global tail-sensitive functional-inequality constants: if polynomial
tails persist, the classical Poincar\'e/LSI conclusions can fail even while local Gaussian
approximation is accurate.  Thus the weighted constants of this section govern that global
behaviour, while the Fisher-information limit \eqref{eq:a2-target} requires the stronger tail
control made explicit in A2 (Section~\ref{sec:a2}, Warning~\ref{warn:a2-tv-fails}). The transport
face of A3 is pursued by A4 (Section~\ref{sec:a4},
Question~\ref{q:a4-modified}) as a
\emph{modified} transport inequality with a cost that is quadratic near $0$ and slower at infinity
--- the $\Wtwo$ analogue of the weighted PI \eqref{eq:a3-weighted-pi}. And the heavy-tailed global
scale of a hierarchical funnel is exactly the Tier-4 obstruction that A5
(Section~\ref{sec:a5}, Question~\ref{q:a5-reparam}) must separate from funnel geometry: a
half-Cauchy in the raw Euclidean scale coordinate $\tau$ defeats Poincar\'e even after coefficient
non-centering.  Passing additionally to $u=\log\tau$ changes the metric by a non-bi-Lipschitz map;
the log-half-Cauchy density is $\operatorname{sech}(u)/\pi$ and has constant $4$, as used in
Proposition~\ref{prop:a3-hierarchical-prior}.
\end{remark}
\end{document}
